\chapter{Conclusion and Future Work}\label{ch:conclusion}

\begin{chapterabstract}
This chapter summarises the results of the project, states their
limitations, and lists directions for future work.
\end{chapterabstract}

\section{Conclusions}

The project studied $k$-dispersion for points in convex position, starting
from the algorithms of Singireddy, Basappa and Mitchell~\citep{sbm}. Its main
outcomes are the following.

\begin{itemize}[leftmargin=*,itemsep=2pt]
\item The bounded search tree of~\citet{sbm} is correct for $k\le4$
(\cref{thm:sbm4}). The proof combines a Delaunay ear argument, which makes the
first step safe for every $k$ (\cref{thm:firststep}), with a lemma on convex
quadrilaterals and a five-point lemma proved by angle chasing. Natural rules
for choosing the starting point fail for $k=5,6,7$ (\cref{prop:rootrules}).
\item The greedy properties behind the fast algorithms for $k=3$ hold for
triples and fail for every $k\ge4$, with explicit small counterexamples
(\cref{sec:greedy}).
\item Under the arc property a single greedy chain is exact
(\cref{thm:arc}); for unimodal point sets this gives an
$O(n\log^2n)$ expected-time algorithm (\cref{thm:arcopt}), compared with the
cubic bounds for general input.
\item A Delaunay certificate for $r$-far sets yields a short correctness proof of
a capped dynamic program that decides the problem in $O(n^3\log k)$ time and
solves it in $O(n^3\log k\log n)$ time (\cref{ch:exact}).
\item All algorithms were implemented with exact arithmetic and agreed with
exhaustive search on every test instance (\cref{ch:results}).
\end{itemize}

\section{Limitations}\label{sec:limitations}

\paragraph{The two-sided search tree.}
$(\mathrm{SBM}_k)$ remains unproved for $k\ge5$, so the bound
$O(2^kn^2\log^2n)$ of~\citet{sbm} is established here only for $k\le4$; the
published argument for general $k$ is a picture argument with three centres of
a Delaunay triangle that we have not been able to complete. Our method does not
extend directly: after a replacement the new point moves away from the
remaining witness points, so the attachment condition of \cref{lem:attach} is
not inherited, and \cref{prop:rootrules} shows that the root must be chosen
with more care than by distances or angles alone. Not even every vertex of a
solution is a good root for $k=4$ (the last instance in
\code{fpt/data/root\_rule\_instances.json}). We conjecture the following
strengthening of $(\mathrm{SBM}_k)$, which holds for $k\le4$ by
\cref{thm:sbm4}: \emph{for every $r$-far $k$-set $T$ there is a vertex of $T$
from which a search that only takes safe steps reaches $k$ points.} The
evidence is experimental: the conjecture held on more than $1.2\cdot10^7$
pairs $(T,r)$ with $k\le7$, including adversarially optimised instances, and
the two-sided search succeeded in all $532{,}833$ random tests and about
$3{,}000$ adversarial runs for $k\le10$. Until a proof is found,
\cref{alg:certified} uses the search only as a sound heuristic.

\paragraph{The arc property.}
The near-linear bound of \cref{thm:arcopt} requires a unimodal point set; for
points exactly on a circle a linear-time algorithm was already
known~\citep{akagi}. For general convex input the arc property must be checked
at every tested radius, which costs $O(n^2)$ per radius, and the adaptive
algorithm may report that it is not applicable. The arc property fails on thin
inputs: in our experiments it held at the optimal radius on every hull of a
random point cloud and on near-circular sets, but never on thin ellipses and
lenses, and applying the decision procedure without the hypothesis gave wrong
answers in $\arcUncheckedWrong\%$ of the cases on thin families. The question
whether contiguous far sets for the points of one solution suffice is open for
$k=4,5$ (\cref{prop:weaker} settles $k=6$).

\paragraph{The exact algorithm.}
The worst-case improvement over~\citet{tw26} is a factor $\log n/\log k$, which
vanishes when $k=n^{\Omega(1)}$, and the per-chord monotonicity behind it is
that of~\citet{tw26}; their algorithm also handles weighted disks of arbitrary
radii, which ours does not. In our experiments the cap made the decision
procedure $\capSpeedupKFour$ times faster for $k=4$ but only $1.1$--$1.4$ times
faster for larger $k$, and the measured growth was $n^{\expCappedFour}$ to
$n^{\expCappedSixteen}$ rather than $n^3$, because the number of binary-search
steps grows with $\log n$ and the table of $\Theta(n^2k)$ entries leaves the
cache. The certified algorithm is no faster than the exact one in the worst
case: when the heuristic bound is not tight, which happened only on thin
inputs, it was about $\certOverhead\%$ slower.

\paragraph{Experiments.}
All experiments use synthetic instances (hulls of random point clouds, points
on or near circles, thin ellipses and lenses) on a single machine. Integer
coordinates are bounded by $2^{29}$ so that the predicates fit into $128$ bits;
unimodal instances with $10^5$ or more points do not fit into this range and
were run with $128$-bit squared distances.



\section{Future work}

\begin{itemize}[leftmargin=*,itemsep=2pt]
\item Prove the claim of~\citet{sbm} for $k\ge5$, or find a counterexample.
The conjectured strengthening stated in \cref{sec:limitations}, that some vertex
of every $r$-far $k$-set is a good root for a search that only takes safe
steps, would suffice.
\item Close the gap between the arc property for all points and for the points
of one solution for $k=4,5$, and find further classes of point sets, beyond
unimodal ones, on which greedy is exact.
\item Remove the remaining $\log n$ factor of the optimisation, for example by
parametric search or by searching directly over the thresholds of the
dynamic program.
\item Extend the capped dynamic program to weighted dispersion and to the
maximum-weight independent set of disks considered by~\citet{tw26}.
\item Evaluate the algorithms on real data, such as candidate facility sites
along a coastline.
\end{itemize}
